% ============================================================================================
% BAB III ANALISIS MASALAH
% Pembagian subbab tidak rigid dan dapat bervariasi. Bab ini minimal berisi analisis kebutuhan
% fungsional, analisis berbagai alternatif solusi yang dapat ditawarkan, dan
% metode pemilihan solusi yang diusulkan.
% ============================================================================================
\cleardoublepage
\chapter{ANALISIS MASALAH}
\label{chap:analisis-masalah}

Bab ini membahas analisis permasalahan yang mendasari pengembangan sistem pembelajaran Algoritma dan Struktur Data (ASD) berbasis \textit{Large Language Model} (LLM) dengan pendekatan \textit{guided learning}. Secara garis besar, bab ini membahas hal-hal berikut:

\begin{enumerate}[1.]
    \item Analisis kondisi pembelajaran ASD saat ini, meliputi fungsi umum sistem yang sudah ada beserta permasalahan yang masih ditemui;
    \item Analisis kebutuhan sistem, mencakup identifikasi masalah pengguna, konsep dasar solusi yang diusulkan, serta rumusan kebutuhan fungsional; dan
    \item Analisis pemilihan solusi, meliputi evaluasi berbagai alternatif teknologi dan penentuan solusi yang paling sesuai untuk mendukung pengembangan sistem.
\end{enumerate}


% ============================================================
\section{Analisis Kondisi Saat Ini}
% ============================================================

Pembelajaran Algoritma dan Struktur Data (ASD) merupakan salah satu komponen penting dalam kurikulum ilmu komputer yang bertujuan membekali mahasiswa dengan kemampuan memahami struktur data, merancang algoritma, serta menganalisis kompleksitas program. Meskipun demikian, berbagai studi menunjukkan bahwa proses pembelajaran ASD masih menghadapi tantangan signifikan, terutama dalam penyampaian konsep yang bersifat abstrak dan dinamis \autocite{Mtaho2024}. Berdasarkan studi literatur dan observasi terhadap praktik pembelajaran saat ini, bagian berikut merangkum fungsi-fungsi utama sistem pembelajaran ASD yang digunakan mahasiswa serta berbagai kendala yang masih ditemui dalam implementasinya.

Adapun fungsi umum yang dimiliki sistem pembelajaran ASD saat ini, yaitu:

\begin{enumerate}
    \item Menyediakan materi pembelajaran dalam bentuk modul, \textit{slide}, dan contoh kode yang dapat digunakan mahasiswa untuk mempelajari konsep dasar, seperti \textit{array}, \textit{linked list}, \textit{stack}, \textit{queue}, \textit{tree}, \textit{graph}, serta algoritma \textit{searching} dan \textit{sorting}.
    \item Menyediakan penjelasan dan ilustrasi visual melalui diagram atau animasi sederhana untuk membantu mahasiswa memahami perubahan struktur data secara bertahap.
    \item Memberikan latihan soal dan kuis yang berfungsi menguji pemahaman mahasiswa terhadap konsep algoritmik tertentu.
    \item Menyediakan forum diskusi atau sesi tanya jawab dengan dosen dan asisten pengajar untuk membantu menjelaskan konsep yang membingungkan atau memperbaiki miskonsepsi.
    \item Menggunakan \textit{environment} pemrograman seperti IDE atau \textit{notebook} untuk memungkinkan mahasiswa mencoba implementasi kode secara langsung.
\end{enumerate}

Meskipun telah menyediakan fitur-fitur tersebut, pendekatan pembelajaran ASD saat ini masih memiliki berbagai permasalahan, antara lain:

\begin{enumerate}
    \item Sistem pembelajaran umumnya bersifat statis dan tidak adaptif terhadap kebutuhan mahasiswa. Materi, contoh, dan penjelasan yang tersedia biasanya disajikan dalam satu format yang sama sehingga tidak dapat menyesuaikan tingkat pemahaman setiap mahasiswa. Kondisi ini membuat mahasiswa yang mengalami kesulitan pada konsep tertentu tidak mendapatkan penjelasan tambahan yang lebih terarah.

    \item Penjelasan mengenai struktur data yang bersifat dinamis, seperti \textit{linked list}, \textit{tree}, atau \textit{graph}, sering kali masih kurang intuitif karena hanya ditampilkan dalam bentuk gambar statis. Hal ini menyulitkan mahasiswa untuk memahami perubahan \textit{pointer}, hubungan antar-\textit{node}, atau proses algoritma secara rinci.

    \item Mahasiswa sering mengandalkan penjelasan dari asisten pengajar atau sumber eksternal seperti \textit{chatbot} LLM umum. Namun, penggunaan LLM secara bebas memiliki keterbatasan, seperti ketidakkonsistenan dengan kurikulum atau potensi terjadinya \textit{hallucination} yang dapat menghasilkan informasi keliru.

    \item Umpan balik terhadap latihan dan pemahaman konsep biasanya tidak diberikan secara cepat. Mahasiswa harus menunggu sesi asistensi atau waktu koreksi, sehingga proses belajar menjadi kurang responsif dan memperlambat perbaikan miskonsepsi.

    \item Sumber referensi materi sering kali tersebar pada berbagai platform seperti LMS, modul digital, forum diskusi, dan sumber eksternal lain, sehingga tidak ada satu wadah terintegrasi yang dapat memandu mahasiswa mengikuti alur belajar secara sistematis.
\end{enumerate}

Saat ini, proses pembelajaran ASD di lingkungan perkuliahan masih sangat bergantung pada metode tradisional yang membutuhkan interaksi langsung dengan dosen atau asisten pengajar. Meskipun LLM modern seperti ChatGPT, Gemini, atau DeepSeek mulai dimanfaatkan oleh mahasiswa, penggunaannya belum terintegrasi dengan kurikulum dan tidak dibatasi oleh struktur pembelajaran tertentu. Kondisi ini membuat mahasiswa berisiko menerima penjelasan yang tidak sesuai urutan materi atau bahkan bertentangan dengan pendekatan yang diajarkan dalam kelas.


% ============================================================
\section{Analisis Kebutuhan}
% ============================================================

Bagian ini menguraikan analisis kebutuhan sistem pembelajaran ASD berbasis \textit{guided learning}. Analisis dimulai dari identifikasi permasalahan yang dihadapi pengguna dalam proses pembelajaran ASD saat ini, kemudian dilanjutkan dengan penggambaran konsep dasar sistem yang diusulkan sebagai solusi, dan diakhiri dengan penjabaran kebutuhan fungsional yang harus dipenuhi oleh sistem.

\subsection{Identifikasi Masalah Pengguna}

Pengguna utama sistem ini adalah mahasiswa yang mempelajari materi Algoritma dan Struktur Data (ASD). Dalam proses pembelajaran saat ini, mahasiswa sering menghadapi berbagai tantangan karena materi ASD memiliki tingkat abstraksi yang tinggi, bersifat dinamis, dan memerlukan kemampuan penalaran algoritmik yang kuat \autocite{Mtaho2024}. Metode pembelajaran tradisional yang mengandalkan modul, \textit{slide}, serta penjelasan langsung dari dosen dan asisten pengajar belum sepenuhnya mampu menjawab kebutuhan mahasiswa dalam memahami konsep secara mendalam dan bertahap. Kondisi ini menyebabkan mahasiswa sering mengalami miskonsepsi pada konsep dasar, kesulitan mengikuti alur materi, serta kurangnya umpan balik cepat ketika mengalami kebingungan.

Selain itu, mahasiswa kerap mencari bantuan melalui \textit{chatbot} LLM umum, namun interaksi yang tidak terarah dan tidak sesuai materi yang diberikan di kelas berpotensi memberikan penjelasan keliru atau tidak konsisten \autocite{White2023PromptPatternCatalog}. Hal ini memperlambat proses belajar dan meningkatkan risiko kesalahan pemahaman. Berikut permasalahan yang dihadapi oleh pengguna tersebut.

\begin{enumerate}
    \item Mahasiswa saat ini belum memiliki sistem pembelajaran yang menyajikan penjelasan konsep ASD secara terstruktur dan bertahap sehingga memudahkan mereka membangun pemahaman dari dasar hingga tingkat yang lebih kompleks.

    \item Mahasiswa tidak memiliki akses terhadap contoh implementasi kode dan ilustrasi langkah demi langkah yang dapat disesuaikan dengan tingkat pemahaman mereka.

    \item Mahasiswa belum memiliki sistem pembelajaran yang mampu memberikan penjelasan, contoh, dan latihan yang konsisten dengan materi karena penggunaan LLM umum tidak dibatasi oleh kurikulum atau konteks materi resmi.

    \item Mahasiswa tidak mendapatkan peringatan atau umpan balik otomatis ketika terjadi miskonsepsi dalam memahami konsep, seperti kesalahan logika pada struktur data atau salah menafsirkan alur algoritma.

    \item Mahasiswa belum memiliki sistem terpadu yang dapat memberikan ringkasan materi yang relevan berdasarkan konteks materi yang sedang mereka pelajari.

    \item Mahasiswa tidak mendapatkan dukungan untuk mengecek pemahaman mereka secara langsung karena latihan dan evaluasi sering kali tidak memberikan penjelasan balik atau pembahasan detail.

    \item Mahasiswa belum memiliki sistem yang mendukung analisis lebih mendalam melalui visualisasi perubahan struktur data dan simulasi algoritma dalam berbagai skenario untuk memperkuat pemahaman konseptual mereka.

    \item Mahasiswa tidak memiliki sarana untuk mengujikan soal dari sumber eksternal (buku, soal ujian terdahulu, atau latihan mandiri) dan mendapatkan evaluasi otomatis terhadap jawaban mereka dalam konteks materi ASD.

    \item Mahasiswa tidak memiliki akses terhadap asisten tanya-jawab yang mampu menjawab pertanyaan kontekstual secara langsung sesuai topik yang sedang dipelajari tanpa harus menunggu sesi asistensi.
\end{enumerate}

Untuk mencari solusi atas masalah-masalah tersebut, perlu disusun kebutuhan fungsional sistem yang diperlukan. Berikut penjabaran mengenai kebutuhan fungsional untuk sistem yang akan dibuat.


\subsection{Konsep Dasar Sistem yang Diusulkan}

Konsep dasar sistem pembelajaran yang diusulkan pada tugas akhir ini berfokus pada pemanfaatan \textit{Large Language Model} (LLM) dalam kerangka \textit{guided learning} yang terstruktur untuk membantu mahasiswa memahami materi Algoritma dan Struktur Data (ASD) secara lebih efektif. Sistem ini dirancang untuk mengatasi berbagai permasalahan pembelajaran yang telah diidentifikasi pada bagian sebelumnya, khususnya terkait sifat materi yang abstrak, perlunya penjelasan bertahap, keterbatasan sumber belajar statis, serta ketidakkonsistenan penjelasan ketika menggunakan \textit{chatbot} umum.

Secara konsep, sistem ini berfungsi sebagai pendamping belajar yang memberikan alur pembelajaran terarah meliputi penjelasan konsep, contoh implementasi, latihan soal, dan ringkasan materi. Setiap tahap pembelajaran dihasilkan oleh LLM namun tetap dikendalikan melalui struktur \textit{prompt} dan konteks materi kuliah yang disediakan oleh dosen. Dengan demikian, sistem tidak hanya menghasilkan penjelasan secara generatif, tetapi melakukannya dengan batasan dan struktur yang selaras dengan kurikulum resmi \autocite{Minaee2025PromptEngineering}.

Selain menyediakan konten terarah, sistem ini juga dirancang untuk memberikan umpan balik otomatis terhadap jawaban mahasiswa pada latihan soal. Dengan adanya umpan balik instan, mahasiswa dapat mengetahui letak kesalahan mereka secara cepat sehingga proses koreksi miskonsepsi dapat berlangsung lebih efektif tanpa harus menunggu sesi asistensi atau evaluasi manual dari pengajar.

Secara keseluruhan, konsep dasar sistem ini adalah menghadirkan platform pembelajaran ASD yang:

\begin{enumerate}
    \item terstruktur melalui tahapan \textit{guided learning} yang konsisten,
    \item adaptif dengan kemampuan LLM menghasilkan penjelasan sesuai konteks materi dan kebutuhan pengguna,
    \item responsif dengan pemberian umpan balik otomatis untuk mendukung koreksi pemahaman, dan
    \item terpadu sebagai satu wadah pembelajaran yang menggabungkan materi, contoh, latihan, dan ringkasan secara sistematis.
\end{enumerate}


\subsection{Kebutuhan Fungsional}

Kebutuhan fungsional atau \textit{Functional Requirements} merupakan kebutuhan-kebutuhan dalam bentuk layanan yang diberikan oleh sistem kepada pengguna. Berikut merupakan tabel kebutuhan fungsional dari sistem pembelajaran Algoritma dan Struktur Data berbasis \textit{Large Language Model} dengan pendekatan \textit{guided learning}.

\begin{longtable}{|L{1.8cm}|L{3cm}|L{7.5cm}|}
\caption{Kebutuhan Fungsional Sistem Pembelajaran ASD Berbasis LLM}
\label{tab:kebutuhan-fungsional}\\
\hline
\textbf{ID} & \textbf{Kebutuhan} & \textbf{Penjelasan} \\
\hline
\endfirsthead

\caption[]{Kebutuhan Fungsional Sistem Pembelajaran ASD Berbasis LLM (lanjutan)}\\
\hline
\textbf{ID} & \textbf{Kebutuhan} & \textbf{Penjelasan} \\
\hline
\endhead

\hline
\multicolumn{3}{l}{\textit{Bersambung ke halaman berikutnya}}\\
\endfoot

\hline
\endlastfoot

FR-01 & Sistem mampu menampilkan alur pembelajaran ASD yang terstruktur &
Sistem dirancang untuk menampilkan alur belajar yang terdiri dari tahapan penjelasan konsep, contoh implementasi, latihan soal, dan ringkasan. Pengguna dapat memilih topik materi, lalu sistem menampilkan urutan langkah yang harus diikuti agar pembelajaran menjadi lebih terarah. \\
\hline

FR-02 & Sistem mampu menghasilkan penjelasan konsep ASD berbasis LLM sesuai konteks materi &
Sistem memanfaatkan LLM dengan struktur \textit{prompt} yang terarah untuk menghasilkan penjelasan konsep yang relevan, akurat, serta konsisten dengan materi yang sedang dipelajari pengguna. \\
\hline

FR-03 & Sistem mampu menampilkan contoh kode dan penjelasan langkah demi langkah &
Sistem dapat menghasilkan contoh implementasi kode untuk topik tertentu, disertai penjelasan langkah demi langkah mengenai alur eksekusi algoritma dan perubahan struktur data. Fitur ini membantu pengguna memahami hubungan antara konsep teoretis dan implementasi program. \\
\hline

FR-04 & Sistem mampu menyediakan latihan soal dan pembahasan otomatis &
Sistem dapat menghasilkan atau menampilkan latihan soal yang terkait dengan materi yang sedang dipelajari pengguna. Setelah pengguna menjawab, sistem memberikan pembahasan otomatis yang menjelaskan alasan jawaban benar maupun salah. \\
\hline

FR-05 & Sistem mampu memberikan umpan balik terhadap miskonsepsi pengguna &
Sistem menganalisis jawaban atau interaksi pengguna untuk mengidentifikasi pola kesalahan umum, kemudian memberikan umpan balik yang menjelaskan letak miskonsepsi serta mengarahkan pengguna kembali ke penjelasan yang sesuai. \\
\hline

FR-06 & Sistem mampu menyajikan ringkasan materi pada akhir sesi belajar &
Sistem dapat menghasilkan ringkasan materi yang telah dipelajari pada suatu sesi, termasuk poin penting, istilah kunci, dan hubungan antar konsep. Ringkasan ini membantu pengguna melakukan \textit{review} dan memperkuat pemahaman konsep. \\
\hline

FR-07 & Sistem mampu menyimpan dan menampilkan riwayat serta progres belajar pengguna &
Sistem menyimpan data riwayat topik yang sudah dipelajari, latihan yang telah dikerjakan, serta status penyelesaian tiap modul. Pengguna dapat melihat progres belajarnya dan melanjutkan pembelajaran dari posisi terakhir dengan mudah. \\
\hline

FR-08 & Sistem mampu menerima dan mengevaluasi soal yang diunggah pengguna &
Sistem menyediakan fasilitas bagi pengguna untuk memasukkan soal dari sumber eksternal beserta jawabannya secara mandiri. Sistem kemudian mengevaluasi jawaban menggunakan LLM berdasarkan konteks materi topik yang sedang dipelajari dan mengembalikan hasil evaluasi berupa skor, komentar, bagian yang sudah benar, bagian yang perlu diperbaiki, serta rekomendasi belajar. \\
\hline

FR-09 & Sistem mampu memberikan penjelasan materi melalui asisten percakapan kontekstual berbasis LLM &
Sistem menyediakan asisten percakapan (\textit{chat asisten materi}) yang memungkinkan pengguna mengajukan pertanyaan seputar materi topik yang sedang dipelajari. Asisten menggunakan konteks materi resmi sebagai landasan jawaban sehingga penjelasan yang dihasilkan tetap relevan dan konsisten dengan kurikulum ASD. Riwayat percakapan dipertahankan selama sesi berlangsung untuk mendukung tanya-jawab yang berkesinambungan. \\
\hline

\end{longtable}

Dengan memenuhi kebutuhan fungsional ini, sistem diharapkan dapat memberikan solusi yang membantu pengguna dalam mengatasi berbagai permasalahan pembelajaran ASD yang telah diidentifikasi sebelumnya.




% ============================================================
\section{Analisis Pemilihan Solusi}
% ============================================================

Bagian ini memaparkan analisis pemilihan teknologi dan pendekatan yang digunakan dalam pengembangan sistem pembelajaran ASD. Analisis mencakup tinjauan terhadap berbagai alternatif solusi yang tersedia, mencakup perbandingan model LLM dan pendekatan pembelajaran, serta justifikasi pemilihan solusi akhir yang paling sesuai dengan kebutuhan sistem.

\subsection{Alternatif Solusi}

Penentuan alternatif solusi dilakukan dengan mempertimbangkan setiap kebutuhan dari permasalahan pembelajaran ASD, khususnya kebutuhan penjelasan yang terarah, konsistensi dengan materi kuliah, serta kemampuan menghasilkan contoh dan latihan secara otomatis. Bagian berikut menguraikan alternatif teknologi yang dapat digunakan dalam pengembangan sistem pembelajaran ASD berbasis LLM, serta pertimbangan pemilihan solusi yang paling sesuai.

\subsubsection{Pemilihan Model LLM}

Terdapat beberapa pilihan model LLM yang populer dalam pengembangan aplikasi edukasi berbasis AI, di antaranya GPT-4o, Gemini 1.5, dan Claude 3. Masing-masing model memiliki kelebihan dan kekurangan dalam hal kemampuan \textit{reasoning}, kualitas penjelasan teknis, serta konsistensi keluaran.

GPT-4o unggul dalam pemahaman konteks yang panjang, penalaran algoritmik, dan kemampuan memberikan penjelasan langkah demi langkah secara stabil, sebagaimana ditunjukkan oleh performanya yang konsisten di atas model-model pendahulunya pada berbagai \textit{benchmark} penalaran dan pemrograman seperti MMLU dan HumanEval \autocite{OpenAI2023GPT4TechnicalReport}. Model ini terbukti efektif untuk menghasilkan contoh kode, menyusun latihan, serta memberikan umpan balik yang terstruktur, menjadikannya kandidat kuat untuk sistem pembelajaran yang memerlukan penjelasan mendalam dalam domain pemrograman \autocite{Jacobs2024LLMFeedbackProgramming}. Selain itu, studi lain menunjukkan bahwa GPT-4 yang berperan sebagai tutor tugas mampu meningkatkan keterlibatan (\textit{engagement}) dan capaian belajar mahasiswa dibandingkan metode belajar konvensional \autocite{Vanzo2024GPT4HomeworkTutor}. Namun, GPT-4o memerlukan biaya penggunaan yang relatif lebih tinggi dibandingkan beberapa alternatif lainnya.

Di sisi lain, Gemini 1.5 menawarkan kapasitas konteks yang sangat besar sehingga cocok untuk sistem yang memerlukan penyertaan modul kuliah atau dokumen panjang dalam \textit{prompt}. Akan tetapi, kualitas \textit{reasoning} Gemini pada permasalahan algoritmik tertentu kadang kurang stabil dan lebih mudah menghasilkan jawaban yang tidak konsisten.

Claude 3 memiliki kemampuan penjelasan natural dan terstruktur, tetapi terkadang kurang presisi dalam contoh kode yang kompleks, khususnya pada topik algoritma tingkat lanjut.

Perbandingan ketiga model LLM tersebut dirangkum pada Tabel~\ref{tab:perbandingan-llm}.

\begin{longtable}{|L{2.9cm}|L{3.3cm}|L{3.3cm}|L{2.4cm}|}
\caption{Perbandingan Model LLM untuk Sistem Pembelajaran ASD}
\label{tab:perbandingan-llm}\\
\hline
\textbf{Kriteria} & \textbf{GPT-4o} & \textbf{Gemini 1.5} & \textbf{Claude 3} \\
\hline
\endfirsthead

\caption[]{Perbandingan Model LLM untuk Sistem Pembelajaran ASD (lanjutan)}\\
\hline
\textbf{Kriteria} & \textbf{GPT-4o} & \textbf{Gemini 1.5} & \textbf{Claude 3} \\
\hline
\endhead

\hline
\multicolumn{4}{l}{\textit{Bersambung ke halaman berikutnya}}\\
\endfoot

\hline
\endlastfoot

\textit{Reasoning} algoritmik & Sangat baik, stabil & Cukup baik, kurang konsisten & Baik, terkadang kurang presisi \\
\hline
Kualitas penjelasan teknis & Mendalam dan sistematis & Cukup baik & Natural, kurang detail pada kode kompleks \\
\hline
Konsistensi keluaran & Tinggi & Sedang & Tinggi \\
\hline
Kapasitas konteks & Besar & Sangat besar & Besar \\
\hline
Biaya penggunaan & Relatif tinggi & Menengah & Menengah \\
\hline
Kesesuaian untuk ASD & Sangat sesuai & Cukup sesuai & Cukup sesuai \\
\hline

\end{longtable}


\subsection{Analisis Penentuan Solusi}

Berdasarkan hasil evaluasi terhadap berbagai alternatif solusi yang telah diuraikan pada subbab sebelumnya, berikut adalah keputusan pemilihan solusi yang digunakan dalam pengembangan sistem pembelajaran ASD berbasis LLM.

\subsubsection{Model LLM yang Dipilih}

Setelah mempertimbangkan seluruh aspek yang relevan, GPT-4o dipilih sebagai model LLM utama yang digunakan dalam sistem ini. Pemilihan GPT-4o didasarkan pada beberapa pertimbangan berikut:

\begin{enumerate}
    \item \textbf{Kemampuan \textit{reasoning} yang tinggi dan stabil.} GPT-4o menunjukkan konsistensi yang baik dalam menghasilkan penjelasan algoritmik yang akurat, terutama untuk topik-topik yang memerlukan penalaran bertahap seperti proses traversal pada \textit{tree} atau analisis kompleksitas algoritma, sejalan dengan hasil pengujian \textit{benchmark} yang menunjukkan keunggulan model ini pada tugas penalaran dan pemrograman \autocite{OpenAI2023GPT4TechnicalReport}.

    \item \textbf{Kualitas contoh kode dan umpan balik yang baik.} Model ini mampu menghasilkan contoh implementasi kode dalam berbagai bahasa pemrograman dengan kualitas yang memadai, disertai penjelasan langkah demi langkah yang mudah dipahami. Efektivitas model keluarga GPT-4 dalam menghasilkan umpan balik pemrograman yang relevan bagi mahasiswa juga telah dikonfirmasi pada penelitian sebelumnya \autocite{Jacobs2024LLMFeedbackProgramming}, sekaligus mendukung perannya sebagai teman belajar digital yang terbukti dapat meningkatkan keterlibatan dan capaian belajar mahasiswa \autocite{Vanzo2024GPT4HomeworkTutor}.

    \item \textbf{Dukungan terhadap \textit{prompt engineering} yang terstruktur.} GPT-4o merespons dengan baik terhadap instruksi dan batasan yang diberikan melalui \textit{system prompt}, sehingga memungkinkan sistem untuk mengendalikan format keluaran sesuai tahapan \textit{guided learning} \autocite{White2023PromptPatternCatalog}.

    \item \textbf{Ekosistem API yang matang.} OpenAI menyediakan API yang stabil dan terdokumentasi dengan baik, mempermudah integrasi dengan \textit{backend} aplikasi.
\end{enumerate}

Dengan pertimbangan tersebut, GPT-4o memberikan keseimbangan terbaik antara kemampuan teknis, stabilitas keluaran, dan kemudahan integrasi untuk mendukung sistem pembelajaran ASD yang terarah, akurat, dan konsisten dengan kurikulum.

\subsubsection{Pendekatan \textit{Guided Learning}}

Selain pemilihan model LLM, pendekatan pembelajaran yang digunakan dalam sistem ini juga merupakan bagian krusial dari solusi yang diusulkan. Pendekatan \textit{guided learning} dipilih karena kemampuannya dalam menyediakan alur belajar yang terstruktur dan terarah, berbeda dengan pendekatan \textit{open-ended chat} yang tidak memiliki batasan konteks. Dengan pendekatan ini, setiap sesi pembelajaran mahasiswa akan mengikuti urutan tahapan yang konsisten, yaitu:

\begin{enumerate}
    \item \textbf{Penjelasan Konsep}: LLM menghasilkan penjelasan teoretis yang disesuaikan dengan topik yang dipilih mahasiswa, mencakup definisi, karakteristik, dan hubungan dengan konsep lain dalam ASD.

    \item \textbf{Contoh Implementasi}: LLM menyajikan contoh kode beserta penjelasan alur eksekusi secara langkah demi langkah, membantu mahasiswa menghubungkan teori dengan praktik pemrograman.

    \item \textbf{Latihan Soal}: Sistem menyediakan soal latihan yang relevan dengan materi, disertai pembahasan otomatis setelah mahasiswa menjawab untuk mendeteksi dan mengoreksi miskonsepsi secara langsung.

    \item \textbf{Ringkasan}: LLM menghasilkan ringkasan poin-poin penting dari sesi belajar yang telah selesai sebagai bahan \textit{review} bagi mahasiswa.
\end{enumerate}

Pendekatan ini memastikan bahwa seluruh konten yang dihasilkan LLM tetap relevan, terstruktur, dan konsisten dengan kurikulum resmi, sekaligus memberikan pengalaman belajar yang lebih terarah dan efektif dibandingkan interaksi bebas dengan \textit{chatbot} umum.
